\documentclass{article}
\usepackage{caption}
\usepackage{tikz}
\usetikzlibrary{backgrounds, positioning, arrows.meta, graphs, trees, graphdrawing}
\usegdlibrary{layered, trees}

\begin{document}
\centering
\centering
\begin{tikzpicture}[scale=1.0, red!100, draw=white!100, node distance=1cm, thick, every node/.style={draw, rectangle}, text height=.7em, text depth=.2em, draw=purple!50, thick, fill=yellow!50]
\begin{tikzpicture}
    % Draw a grid for reference
    \draw [help lines] (0,0) grid (7,7);

    % Define two layered graphs
    \graph [green!100] [layered layout, grow=down] {
        i [y=7.5, x=-0.5] -> {1 -> 2 -> 3 -> 4 -> 5 -> 6 -> 7};
    };
    \graph [layered layout, grow=down] {
        j [x=3.5, y=8.5] -> {a, b, c, d, e, f, g};
    };

    % Add a label at the specified coordinates
    \draw (3.5, 9.5) node {The domain and system layer are defined iteratively through generative adjunctions};
\end{tikzpicture}

\end{tikzpicture}
\captionof{figure}{}
\vspace{5ex}

\end{document}
